% Section content template for individual Educative lessons
% This file contains the actual content for a specific section
% Generated from Educative API section content

% Section: Types of UML Diagrams
% Section ID: 4722803754663936
% Section Slug: types-of-uml-diagrams
% Generated: 2025-12-08T22:21:10.413702

UML diagrams can be classified into two categories:

\begin{itemize}
\item
Structural UML diagrams
\item
Behavioral UML diagrams
\end{itemize}
\\
The following illustration represents the subcategories of structural and behavioral UML diagrams:

\begin{center}
\fbox{\begin{minipage}{0.9\textwidth}
\centering
\vspace{0.2cm}
\textbf{\large Types of UML diagrams}
\\[0.2cm]
\textit{[Interactive Mind Map - See online version for interactive visualization]}
\vspace{0.2cm}
\end{minipage}}
\end{center}

\vspace{0.3cm}
\begin{itemize}
\item \textbf{UML Diagrams}
\begin{itemize}
\item \textbf{Structural Diagrams}
\begin{itemize}
\item Package Diagram
\item Class Diagram
\item Object Diagram
\item Composite Structure Diagram
\end{itemize}
\begin{itemize}
\item \textbf{Implementation Diagrams}
\begin{itemize}
\item Component Diagram
\item Deployment Diagram
\end{itemize}
\end{itemize}
\end{itemize}
\begin{itemize}
\item \textbf{Behavioral Diagrams}
\begin{itemize}
\item Use Case Diagram
\item Activity Diagram
\item State Machine Diagram
\end{itemize}
\begin{itemize}
\item \textbf{Interaction Diagrams}
\begin{itemize}
\item Sequence Diagram
\item Communication Diagram
\item Interaction Overview Diagram
\item Timing Diagram
\end{itemize}
\end{itemize}
\end{itemize}
\end{itemize}

\subsection{Structural UML diagrams}\label{structural-uml-diagrams}
\\
Structural diagramsrepresent the static structure of the system. They never depict the system's dynamic behavior. The most commonly used structural diagram in software development is the \emph{class diagram}.

\subsection{Behavioral UML diagrams}\label{behavioral-uml-diagrams}
\\
Behavioral diagrams represent the dynamic behavior of elements in the system. All systems experience dynamic occurrences. In object-oriented programming, we use behavioral diagrams to illustrate the dynamic behavioral semantics of a problem or its implementation. The most commonly used behavioral diagrams are \emph{use case diagrams}, \emph{activity diagrams}, and \emph{sequence diagrams}.

\subsubsection{UML diagrams used in the OOD interview}\label{uml-diagrams-used-in-the-ood-interview}
\\
In a typical OOD interview process, the following are the most widely asked UML diagrams:

\begin{itemize}
\item
Use case diagram
\item
Class diagram
\item
Sequence diagram
\item
Activity diagram
\end{itemize}
\begin{quote}
\textbf{Note:} The class diagram belongs to the structural diagrams and the use case, sequence, and activity diagrams belong to the behavioral diagrams category of the UML.
\end{quote}
\\
The next few lessons will cover all the four UML diagrams above, in detail.

% End of section content